% TOP_PROBLEM: {"id":30007696,"problem_number":"LOCAL-30007696","title":"Organizational complements","category_id":21,"category":{"id":21,"name":"economics","display_name":"Economics","description":"Open economic theory statements, published research agendas, and proposed scoped specifications; question provenance and historical priority are qualified per record.","slug":"economics","order_index":21,"created_at":"2026-10-08T00:00:00Z"},"status":"provenance_pending","external_url":null,"legacy_ids":[],"published":false,"statement":"Choose task allocation, verification workflows and decision authority that complement AI-assisted R&D in a team-level factorial design.","economics":{"original_id":185,"field":"Innovation and AI economics","kind":"design","formalization_basis":"proposed_specification","source_status":"not_verified","priority_status":"not_established","priority_note":"No exact open-problem passage was verified. The supporting publication does not establish priority or unresolved status.","earliest_verified_reference":null,"current_reference":{"authors":["Iain M. Cockburn","Rebecca Henderson","Scott Stern"],"title":"The Impact of Artificial Intelligence on Innovation","year":2018,"url":"https://www.nber.org/papers/w24449","doi":"10.3386/w24449","locator":"Primary publication abstract; exact organizational-redesign agenda paragraph not retrieved","evidence_summary":"Motivates AI as a method of invention and R&D reorganization, but does not verify the specific team/workflow design question."},"formalization_note":"The cited primary work supports the subject or supplies solutions in particular settings, but no exact open-question passage for this specification was verified. The mathematical target and restrictions are proposed here, with no canonical-conjecture or priority claim."},"statusQualification":"Provenance pending: an explicit published statement of this exact open question was not verified. This is a catalogue-authored candidate specification, not a certified open conjecture. The cited primary work supports the subject or supplies solutions in particular settings, but no exact open-question passage for this specification was verified. The mathematical target and restrictions are proposed here, with no canonical-conjecture or priority claim.","statusReviewedAt":"2026-10-08"}
% FIRST_POSTED: 2026-10-08
% ENTRY_KIND: statement_only
% !TeX program = lualatex
\documentclass[11pt]{article}
\usepackage{fontspec}
\usepackage[a4paper,margin=25mm]{geometry}
\usepackage{amsmath,amssymb,mathtools}
\usepackage{xurl}
\usepackage[hidelinks]{hyperref}
\newcommand{\sref}[2]{\hyperlink{source:#1:#2}{[S#2]}}
\setlength{\emergencystretch}{3em}
\begin{document}
% problemId:30007696
\section*{Organizational complements}\label{30007696}
\subsection{Definitions and mathematical statement}
Consider research teams inside firms conducting comparable product-development projects, with a fixed two-year validation horizon. Let \(A\in\{0,1\}\) indicate access to an AI research assistant, and let \(Z\in\mathcal Z\) denote one of finitely many preregistered organizational designs. Each design specifies task allocation, experiment approval rules, independent verification responsibilities and the distribution of decision authority. Design \(z_0\) is the existing organization. Let \(Y_i(a,z)\) be independently validated discovery value and \(C_i(a,z)>0\) the real labor, computation, experimentation and review cost for team \(i\). Define \(P(a,z)=E[Y_i(a,z)]/E[C_i(a,z)]\) over eligible research teams, and the organizational complementarity
\[
\Delta(z)=P(1,z)-P(0,z)-P(1,z_0)+P(0,z_0).
\]
Randomize AI access and organizational designs across sufficiently isolated teams, or across connected team clusters where knowledge sharing occurs. Hold project eligibility and evaluation criteria fixed, include failed projects, and distinguish faster drafting from reproducible technical discoveries. The design problem is to choose \(z^\ast\in\arg\max_{z\in\mathcal Z}P(1,z)\), subject to stated safety and review-capacity constraints. Which redistribution of research work and authority creates positive complementarity with AI, and what mechanisms explain it? This specification concerns AI-assisted R\&D organization rather than economy-wide adoption delays or aggregate intangible investment.

\par\textbf{Formulation and scope.} The cited primary work supports the subject or supplies solutions in particular settings, but no exact open-question passage for this specification was verified. The mathematical target and restrictions are proposed here, with no canonical-conjecture or priority claim.

\par\textbf{Open-question provenance.} An explicit published open-question statement was not verified. Sources below are supporting literature and must not be treated as a first listing.

\par\textbf{Historical priority.} No exact open-problem passage was verified. The supporting publication does not establish priority or unresolved status.

\subsection{Short English statement}
Choose task allocation, verification workflows and decision authority that complement AI-assisted R\&D in a team-level factorial design.

\subsection{Sources}
\begin{itemize}\raggedright
\item[\textnormal{[S1]}]\hypertarget{source:30007696:1}{} Iain M. Cockburn, Rebecca Henderson, Scott Stern. 2018. The Impact of Artificial Intelligence on Innovation. Primary publication abstract; exact organizational-redesign agenda paragraph not retrieved.
\url{https://www.nber.org/papers/w24449}.
DOI: 10.3386/w24449.
{\small\textit{Supporting or current literature; see evidence qualification. Motivates AI as a method of invention and R\&D reorganization, but does not verify the specific team/workflow design question.}}
\end{itemize}
\end{document}
